\scanpage{95}\mspage{47}
B)
% slanted
\margin{On suppose maintenant $2\cdot1_\Lambda\ne0$}
\[
\left\{\begin{aligned}
G'&=\mathrm{Sl}(2,\Lambda)'=\deleted{\unsure{S}}G/\pm1\\
H'&=\text{image dans }G'\text{ du groupe }H\text{ de (49)}\\
&\hspace{3em}\text{par }p:G\to G'\quad\text{(can)}
\end{aligned}\right.\tag*{(51)}
\]
Soit $g'$ un élément de $G'$, provenant de
\[
g=\begin{pmatrix}a&b\\c&d\end{pmatrix},\qquad ad-bc=1
\]
et considérons \deleted{\unsure{pour $t$|pour le}} \deleted{les}
\[
\operatorname{int}(g')\cdot\underbrace{p\begin{pmatrix}1&t\\0&1\end{pmatrix}}_{h'=p(h)}
=p\ \operatorname{int}(g)\begin{pmatrix}1&t\\0&1\end{pmatrix}
=p\begin{pmatrix}1-act&a^2t\\-c^2t&1+act\end{pmatrix}
\]
% slanted
\margin{\uline{NB} l'image inv. de $H'$ dans $G$ est formé des $\left(\begin{smallmatrix}\pm1&\bullet\\0&\bullet\end{smallmatrix}\right)$}
on a
\[
\operatorname{int}(g')(h')\in H'\Longleftrightarrow
\left|\begin{aligned}-c^2t&=0\\1-act&\in\pm1\end{aligned}\right.\tag*{(52)}
\]
Donc
\[
g'\in N_{G'}(H')\Longleftrightarrow
\left\{\begin{aligned}&c^2=0\\&\forall t\in\Lambda,\ \text{on a}\quad act=0\text{ ou }act=2\end{aligned}\right.
\]

Considérons le cas $t=1$, on aura donc
\[
ac=0\quad\text{ou}\quad ac=2
\]
dans le premier cas, on a donc (puisque
$c^2=0\Longrightarrow a\in\Lambda^*$) $c=0$, ce qui exclut d'ailleurs
$ac=2$ (puisque $2_\Lambda\ne0$ par hypothèse). Dans
le deuxième cas, on doit avoir, en élevant
au carré
\[
4_\Lambda=0
\]
et de plus, la \inserted{$2^{\text e}$} relation (52) \deleted{\unsure{implique}} \inserted{s'écrit}
\[
1-2t\in\pm1\qquad\forall t\in\Lambda\qquad\text{i.e.}
\]
\[
\forall t\in\Lambda,\ \text{on a}\quad 2t=0\quad\text{ou}\quad 2t=2\tag*{(53)}
\]
ce qui implique que l'idéal $2\Lambda$ est réduit à deux
% archive sheet number: 94
